% Fixed132 substantive insertion.
% Crossed quadratic-etale multiplication, exact wild adjoint blocks,
% 2-primary return parity, Schreier exchange Laplacian and
% formal elimination of the second wild off-isotypic matrix coordinate.

\subsection{Crossed quadratic-\'etale wild transport, return parity, and
Schreier first-exchange syzygies}
\label{subsec:dyadic-f132-crossed-schreier-exchange}

The curved two-block carrier of
Theorem~\ref{thm:dyadic-f131-curved-tame-transport}
admits a stronger \emph{ring-level} calculation.
The integral native frame $(I,X,Y,XY)$ gives an
actual quadratic-\'etale crossed-product algebra,
rather than an arbitrary block decomposition of
an abstract four-dimensional module.  This structure
simultaneously controls the nonlinear wild matrices,
their inverse and adjoint Fox coefficients, the exact
order of their two-step diagonal return defects,
and the cyclic first-exchange relations across the
three Schreier cosets.  The marked wild relation
then allows one of the two off-isotypic wild
coordinates to be eliminated formally.
None of these statements assumes a nonexistent
global Galois-equivariant splitting.

Keep the full framed residual $S_3$ family of fixed131.
For a complete local deformation algebra $A$ put
\[
X=\rho_A(\tau),\quad Y=\rho_A(\sigma),\quad
C=A[X]\simeq A[T]/(T^2+T+1).
\]
Let $\iota:C\to C$ be the nontrivial
$A$-automorphism $\iota(X)=X^2$ and write
\begin{equation}\label{eq:dyadic-f132-crossed-setup}
q=Y^2\in A^\times,\qquad
N(a)=a\iota(a)\in A,\quad
\operatorname{Tr}(a)=a+\iota(a)\in A.
\end{equation}
The exact tame relation gives
$Y a=\iota(a)Y$ for every $a\in C$,
and Theorem~\ref{thm:dyadic-f131-native-tame-four-frame}
gives the canonical decomposition
$M_2(A)=C\oplus CY$.
Below the letter $a$ denotes a
coefficient in $C$, not the
fixed130 generator named $a$.

\begin{theorem}[A complete integral crossed-\'etale matrix calculus]
\label{thm:dyadic-f132-crossed-etale-multiplication}
For $a,b,c,d\in C$ the multiplication
in $M_2(A)$, expressed in the integral
native frame, is
\begin{equation}\label{eq:dyadic-f132-crossed-product-law}
\boxed{
(a+bY)(c+dY)
=
\bigl(ac+q\,b\iota(d)\bigr)
+\bigl(ad+b\iota(c)\bigr)Y.
}
\end{equation}
Its determinant and adjugate are
\begin{equation}\label{eq:dyadic-f132-crossed-determinant}
\boxed{
\det_A(a+bY)=N(a)-qN(b)=:\Delta(a,b),
\qquad
(a+bY)^\natural=\iota(a)-bY.
}
\end{equation}
In particular, when $\Delta(a,b)\in A^\times$,
\begin{equation}\label{eq:dyadic-f132-crossed-inverse}
\boxed{\qquad
(a+bY)^{-1}
=\frac{\iota(a)-bY}{N(a)-qN(b)}.
\qquad}
\end{equation}
The wild matrices
\[
W_j=\rho_A(x_j)=a_j+b_jY,\qquad
a_j\equiv1\pmod{\mathfrak m C},\quad
b_j\in\mathfrak m C,
\quad j=0,1,
\]
have $a_j\in C^\times$ and
$\Delta_j=N(a_j)-qN(b_j)\in A^\times$.
The difference between centralizer
and noncentralizer transport is
the exact formula
\begin{equation}\label{eq:dyadic-f132-wild-commutator}
\boxed{\qquad
[W_j,X]_{\rm ring}
=b_j(X^2-X)Y,\qquad
(X^2-X)^2=-3I_2\in C^\times.
\qquad}
\end{equation}
Consequently $W_j$ commutes with
the canonical tame generator $X$
if and only if $b_j=0$.
The full matrix multiplication,
inversion and centralizer tests
are thus finite integral operations,
with no divisions other than the
units $3$ and $\Delta(a,b)$.
\end{theorem}

\begin{proof}
The identity $Ya=\iota(a)Y$
and $Y^2=q$ immediately give
\eqref{eq:dyadic-f132-crossed-product-law}.
Multiplying $a+bY$ on either side by
$\iota(a)-bY$ cancels its
off-isotypic part and leaves
$N(a)-qN(b)$, proving the
inverse identity whenever
the scalar is a unit.
To identify that scalar with
the ordinary determinant,
make the faithfully flat
\'etale base change splitting
$C\simeq A\times A$.
In that basis, $a$ is diagonal,
$Y$ exchanges the two factors
up to the scalar $q$, and
the displayed scalar is
the ordinary $2\times2$
determinant.  This identity
descends to $A$.
The residual wild values are
identity, so their $a_j$
and $\Delta_j$ are units.

As $Xa=aX$ for $a\in C$
and $YX=X^2Y$, one has
\[
(a+bY)X-X(a+bY)
=b(X^2-X)Y.
\]
Finally $X^2=-1-X$ implies
$(X^2-X)^2=(-1-2X)^2=-3$.
The scalar $-3$ is a $2$-adic
unit, so $X^2-X$ is invertible
and the centralizer condition
is equivalent to $b_j=0$.
\end{proof}

For an arbitrary
$W=a+bY\in\mathrm{GL}_2(A)$
set $\Delta=N(a)-qN(b)$.
The universal tame projector
gives $P_0:M_2(A)\to C$
and $P_1:M_2(A)\to CY$,
and the adjoint transport
$\mathcal R_W(Z)=WZW^{-1}$
has four actual finite
$A$-linear blocks
\[
\mathcal R_W=
\begin{pmatrix}
\mathcal A_W&\mathcal B_W\\
\mathcal C_W&\mathcal D_W
\end{pmatrix}
\quad\text{on }C\oplus CY.
\]

\begin{theorem}[Closed twisted adjoint blocks and the extra dyadic return factor]
\label{thm:dyadic-f132-adjoint-return-parity}
With this notation, for $z,d\in C$
the four adjoint blocks are
\begin{equation}\label{eq:dyadic-f132-adjoint-four-blocks}
\boxed{
\begin{aligned}
\mathcal A_W(z)
&=\frac{N(a)z-qN(b)\,\iota(z)}{\Delta},\\
\mathcal C_W(z)
&=\frac{ab}{\Delta}\,(\iota(z)-z)\,Y,\\
\mathcal B_W(dY)
&=\frac{q}{\Delta}
\left(b\iota(a)\iota(d)
       -a\iota(b)d\right),\\
\mathcal D_W(dY)
&=\frac{a^2d-qb^2\iota(d)}{\Delta}\,Y.
\end{aligned}}
\end{equation}
In particular, for a wild
matrix $W=a+bY$, the
truncated invariant transport
satisfies the \emph{exact}
quadratic deviation formula
\begin{equation}\label{eq:dyadic-f132-invariant-diagonal-drift}
\boxed{\qquad
\mathcal A_W(z)-z
=\frac{qN(b)}{\Delta}\,
(z-\iota(z)).
\qquad}
\end{equation}
Thus the first-order
off-isotypic parameter $b$
cannot affect the invariant
diagonal carrier linearly.

For any two matrices
$G=a_g+b_gY$, $H=a_h+b_hY$
of invertible determinant
$\Delta_g,\Delta_h$,
the two projected
two-step return defects
\[
\mathscr F_{00}(G,H)
=\mathcal A_{GH}
  -\mathcal A_G\mathcal A_H
=\mathcal B_G\mathcal C_H,
\quad
\mathscr F_{11}(G,H)
=\mathcal D_{GH}
  -\mathcal D_G\mathcal D_H
=\mathcal C_G\mathcal B_H
\]
are given \emph{exactly} by
\begin{equation}\label{eq:dyadic-f132-return-zero}
\boxed{
\begin{aligned}
\mathscr F_{00}(G,H)(z)
={}&\frac{q}{\Delta_g\Delta_h}
\Bigl(
 b_g\iota(a_g)\iota(a_h)\iota(b_h)
 +a_g\iota(b_g)a_hb_h
\Bigr)(z-\iota(z)),
\\[4pt]
\mathscr F_{11}(G,H)(dY)
={}&\frac{2q\,a_gb_g}
{\Delta_g\Delta_h}
\Bigl(
a_h\iota(b_h)d-
b_h\iota(a_h)\iota(d)
\Bigr)Y.
\end{aligned}}
\end{equation}
Let $\mathfrak J$ be the wild
off-isotypic ideal in the
complete residual deformation
ring, generated by the two
$C$-coordinates $b_0,b_1$.
For every $g,h\in\Gamma$,
the off-diagonal blocks of
$\mathcal R_g,\mathcal R_h$
have coefficients in
$\mathfrak J$, and hence
\begin{equation}\label{eq:dyadic-f132-return-filtration}
\boxed{\quad
\mathscr F_{00}(g,h)\in
\mathfrak J^2
\operatorname{Hom}_A(C,C),
\qquad
\mathscr F_{11}(g,h)\in
2\mathfrak J^2
\operatorname{Hom}_A(CY,CY).
\quad}
\end{equation}
Consequently on \emph{any}
characteristic-two base
change, $\mathcal D_{gh}
=\mathcal D_g\mathcal D_h$
is a literal equality of
the projected second-diagonal
block, even when the
off-diagonal wild transport
is nonzero.
No analogous statement
is made for $\mathcal A$;
its first possible failure
is genuinely quadratic.
\end{theorem}

\begin{proof}
Insert the inverse formula
\eqref{eq:dyadic-f132-crossed-inverse}
into $WzW^{-1}$ and
$W(dY)W^{-1}$ and apply
\eqref{eq:dyadic-f132-crossed-product-law}.
The central and $Y$ parts
give exactly
\eqref{eq:dyadic-f132-adjoint-four-blocks}.
Since $\Delta=N(a)-qN(b)$,
the invariant-block drift
simplifies to
\eqref{eq:dyadic-f132-invariant-diagonal-drift}.

Block multiplication for
$\mathcal R_{GH}
=\mathcal R_G\mathcal R_H$
gives
$\mathscr F_{00}=\mathcal B_G\mathcal C_H$
and
$\mathscr F_{11}=\mathcal C_G\mathcal B_H$.
Substituting the four
closed blocks and using
\[
\iota(\iota(z)-z)=z-\iota(z)
\]
gives the first displayed
return formula.
For the second, write
\[
\mathcal B_H(dY)
=\frac{q}{\Delta_h}
\bigl(b_h\iota(a_h)\iota(d)
       -a_h\iota(b_h)d\bigr).
\]
The difference of its
Galois conjugate and itself is
\[
\iota(\mathcal B_H(dY))
-\mathcal B_H(dY)
=\frac{2q}{\Delta_h}
\bigl(a_h\iota(b_h)d
      -b_h\iota(a_h)\iota(d)\bigr).
\]
Applying $\mathcal C_G$
then produces the
extra factor two in
\eqref{eq:dyadic-f132-return-zero}.
The formula is integral
even before inverting two.

Modulo the closed ideal
$\mathfrak J$, both wild
transport matrices lie in
$C$, while $\sigma$ acts
through the \'etale involution
and $\tau$ through $C$.
Thus every group word
normalizes $C$ modulo
$\mathfrak J$, and its
projected off-diagonal
adjoint blocks vanish
modulo $\mathfrak J$.
Their products therefore
lie in $\mathfrak J^2$.
The explicit second
return formula strengthens
the second inclusion to
$2\mathfrak J^2$.
After characteristic-two
base change the latter
vanishes identically.
The first return may remain
nonzero, as the next
Schreier calculation shows.
\end{proof}

\begin{theorem}[Exact three-coset Schreier exchange and zero-sum curvature]
\label{thm:dyadic-f132-schreier-exchange}
For each wild source generator
$x_j$, write
$W_j=a_j+b_jY$ in
the native quadratic-\'etale
frame.  Its three
actual Schreier conjugates
from fixed128 are
\[
w_{ij}=\tau^i x_j\tau^{-i},
\qquad i=0,1,2.
\]
Their transport matrices
are given \emph{exactly} by
\begin{equation}\label{eq:dyadic-f132-schreier-wild-packet}
\boxed{\quad
\rho_A(w_{ij})
=a_j+b_jX^{2i}Y,
\qquad
b_{ij}:=b_jX^{2i}\in C.
\quad}
\end{equation}
Therefore each fixed-$j$
Schreier packet has the
precise first-exchange
identity
\begin{equation}\label{eq:dyadic-f132-schreier-zero-sum}
\boxed{\qquad
b_{0j}+b_{1j}+b_{2j}
=b_j(1+X^2+X)=0.
\qquad}
\end{equation}
Moreover
$N(b_{ij})=N(b_j)$,
so every conjugate has
the same determinant
$\Delta_j=N(a_j)-qN(b_j)$.

For one fixed $j$ define
the $3\times3$ family of
quadratic return maps
\[
\mathscr K_{i\ell}(z)
:=\mathscr F_{00}
(\rho_A(w_{ij}),
 \rho_A(w_{\ell j}))(z),
\qquad i,\ell=0,1,2.
\]
They have the explicit
\'etale-valued formula
\begin{equation}\label{eq:dyadic-f132-schreier-quadratic-array}
\boxed{
\begin{aligned}
\mathscr K_{i\ell}(z)
=\frac{qN(b_j)}{\Delta_j^2}
\Bigl(
 \iota(a_j)^2X^{2i+\ell}
 +a_j^2X^{i+2\ell}
\Bigr)\,(z-\iota(z)).
\end{aligned}}
\end{equation}
In particular \emph{every}
row and column sum
of this quadratic
exchange array
is exactly zero:
\begin{equation}\label{eq:dyadic-f132-schreier-curvature-zero-sums}
\boxed{\quad
\sum_{i=0}^2\mathscr K_{i\ell}=0,
\qquad
\sum_{\ell=0}^2\mathscr K_{i\ell}=0.
\quad}
\end{equation}
The second-diagonal
return array satisfies
the same zero-sum
identities and, by
\eqref{eq:dyadic-f132-return-filtration},
is divisible by two.

In the important
normalized case
$a_j=1$, the complete
first-diagonal
exchange array becomes
a scalar multiple of
the three-vertex
complete-graph Laplacian:
\begin{equation}\label{eq:dyadic-f132-schreier-k3-laplacian}
\boxed{
\mathscr K_{i\ell}(z)
=
\frac{qN(b_j)}
{(1-qN(b_j))^2}
(3\delta_{i\ell}-1)
(z-\iota(z)).
}
\end{equation}
Thus a genuine
quadratic return
profile may survive
on the invariant
diagonal block
even though the
full Galois transport
is flat and the
degree-one Schreier
exchange sum is zero.

Equations
\eqref{eq:dyadic-f132-schreier-zero-sum}
and
\eqref{eq:dyadic-f132-schreier-curvature-zero-sums}
are finite
\emph{coefficient-level}
Schreier exchange
identities.  They
are not asserted to
be defining group
relators or an
integral replacement
for the full Fox
relation-syzygy module.
\end{theorem}

\begin{proof}
The exact tame
relation implies
$XYX^{-1}=X^2Y$.
Since the $a_j$
and $b_j$ commute
with $X$,
\[
X^iW_jX^{-i}
=a_j+b_jX^{2i}Y.
\]
The equality
$1+X+X^2=0$
proves the first
packet identity.
Also $N(X)=X\iota(X)=X^3=1$,
hence every Schreier
conjugate has the same
norm and determinant.

Substitute
$a_g=a_h=a_j$,
$b_g=b_jX^{2i}$ and
$b_h=b_jX^{2\ell}$
into the first line of
\eqref{eq:dyadic-f132-return-zero}.
Using
\[
b_jX^{2i}\,\iota(b_jX^{2\ell})
=N(b_j)X^{2i+\ell},
\quad
\iota(b_jX^{2i})\,b_jX^{2\ell}
=N(b_j)X^{i+2\ell}
\]
gives the
explicit quadratic array.
Each row and column
sum vanishes by
$1+X+X^2=0$.
The second array
is linear in the
first off coefficient
$b_jX^{2i}$ and
linear in the
second off coefficient
$b_jX^{2\ell}$
or its conjugate,
so the same zero-sum
argument applies;
its extra factor
two is
Theorem~\ref{thm:dyadic-f132-adjoint-return-parity}.

If $a_j=1$, its
quadratic coefficient
is
\[
X^{2i+\ell}+X^{i+2\ell}
=
\begin{cases}
2,&i=\ell,\\
-1,&i\ne\ell.
\end{cases}
\]
Indeed for $i\ne\ell$
the two exponents
are congruent to
$1$ and $2$ modulo
three.  Substitution
yields the complete-graph
Laplacian formula.
\end{proof}

The fixed131 split
ideal was originally
defined by the four
off-isotypic scalars of
$b_0$ and $b_1$.
The next result uses
the actual Roe--Turturean
wild relation, rather
than merely the crossed
algebra, to show that
two of these four
generators are formally
redundant.

\begin{theorem}[Formal elimination of the second wild off-isotypic variable]
\label{thm:dyadic-f132-formal-wild-elimination}
Work on the complete
framed $S_3$ residual
deformation chart.
Write
\[
b_j=\eta_j+\xi_jX\in C,
\quad
a_j=1+\text{higher terms in }C,
\qquad j=0,1.
\]
Let $F_{\rm wild}$ denote
the \emph{literal}
Roe--Turturean marked
wild relation matrix
minus the identity,
and let
\[
F_1=(1-\Pi)F_{\rm wild}
\in CY
\]
be its two
off-isotypic scalar
equations.
After imposing the
tame relation, and
treating all diagonal
wild coordinates
and $b_0$ as
parameters, the
Jacobian of $F_1$
with respect to
the two coordinates
$(\eta_1,\xi_1)$
is invertible
at the residual
$S_3$ point:
\begin{equation}\label{eq:dyadic-f132-wild-off-jacobian-unit}
\boxed{\quad
D_{b_1}F_1\big|_{\bar\rho}
=\operatorname{Ad}_{\bar S}^{-1}
\big|_{\bar M_1},
\qquad
\det D_{b_1}F_1\big|_{\bar\rho}
\in\mathbf F_2^\times.
\quad}
\end{equation}
There is therefore
a \emph{unique}
formal solution
\begin{equation}\label{eq:dyadic-f132-wild-b1-implicit-series}
\boxed{\quad
b_1=\mathcal G(b_0;\,
X,Y,a_0,a_1)
\quad}
\end{equation}
of the two off
relation equations
over the complete
tame/diagonal base.
It is functorial
under complete
local coefficient
base change and
has the exact
properties
\begin{equation}\label{eq:dyadic-f132-wild-b1-order}
\boxed{\quad
\mathcal G(0;\cdots)=0,
\qquad
b_1\in\mathfrak m_R(b_0)
\quad\text{in }R_{\bar\rho}^{\square}.
\quad}
\end{equation}
Here $(b_0)$ means
the ideal generated
by its two scalar
coordinates
$\eta_0,\xi_0$,
not the principal
$C$-module ideal
mistaken for one
$A$-valued equation.

Consequently the
entire tame-isotypic
split ideal has
the sharper exact
description
\begin{equation}\label{eq:dyadic-f132-two-generator-split-ideal}
\boxed{\qquad
\mathfrak I_{\rm split}
=(\eta_0,\xi_0)
\subset R_{\bar\rho}^{\square},
\qquad
\eta_1,\xi_1\in
\mathfrak m_R
(\eta_0,\xi_0).
\qquad}
\end{equation}
The two generators
are minimal at the
residual point by
Corollary~\ref{cor:dyadic-f131-split-tangent-codimension}.
Thus \emph{two}, not
four, independent
equations suffice
scheme-theoretically
on the complete
framed residual chart.

This formal
elimination does
not prove that
$\mathfrak I_{\rm split}$
is a regular sequence,
that its quotient
is smooth, or that
a canonical minimal
Demu\v{s}kin Fox
relation module has
been constructed.
It removes only the
off-isotypic
second-wild-variable
redundancy in the
genuine completed
Roe--Turturean
relation locus.
\end{theorem}

\begin{proof}
The native matrix
frame of
Theorem~\ref{thm:dyadic-f131-native-tame-four-frame}
provides a unit
formal change of
coordinates from the
four entries of each
wild matrix $W_j$
to the two
$C$-coordinates
of $a_j$ and the
two of $b_j$.
After imposing the
exact tame relation,
the two $CY$
components of the
single wild matrix
relator are formal
power series in
these coordinates.

The coefficient
derivative on the
residual $M_1$
natural summand is
the wild row
of
Theorem~\ref{thm:dyadic-s3-adjoint-analytic-smith}:
\[
d^1_{\rm wild}|_{\bar M_1}
(c,d)=-2c+\bar S^{-1}d.
\]
The first summand
vanishes in residue
characteristic two,
and $\bar S^{-1}$
is an invertible
linear transformation
of the rank-two
off-isotypic module.
This is exactly
\eqref{eq:dyadic-f132-wild-off-jacobian-unit}.

Apply the formal
implicit function
theorem over the
complete local
tame/diagonal
base to the two
off relation equations.
The invertible
Jacobian in the
two $b_1$
coordinates yields
a unique compatible
solution $b_1=\mathcal G$.
If $b_0=b_1=0$,
all wild generator
matrices belong to
$C$, while the
tame generators
normalize $C$.
Every marked
auxiliary word,
including
$\omega_2$ powers,
then lies in
$C$ or its
normalizer as
appropriate;
the complete
wild relation
evaluates in
$C$, so its
off component
vanishes
\emph{identically}
for arbitrary
diagonal coordinates.
By uniqueness of
the formal solution,
$\mathcal G(0;\cdots)=0$.
Hence every
coefficient of
$b_1$ belongs
to the two-generated
ideal $(\eta_0,\xi_0)$.

Moreover the
residual differential
with respect to
$b_0$ is zero,
as the displayed
wild row shows.
Thus the linear
part of $\mathcal G$
in $(\eta_0,\xi_0)$
has coefficients
in the maximal
ideal, and its
higher terms also
lie in
$\mathfrak m_R
(\eta_0,\xi_0)$.
This proves
\eqref{eq:dyadic-f132-wild-b1-order}.
The split ideal
was the combined
four-coordinate
ideal $(b_0,b_1)$
in fixed131.
Since $b_1$ now
lies in $(b_0)$,
its exact equality
with $(\eta_0,\xi_0)$
follows.
The two-dimensional
split tangent
codimension in
fixed131 shows that
neither of the
two remaining
generators can be
discarded.
\end{proof}

\begin{corollary}[One-source Schreier control of all wild exchange defects]
\label{cor:dyadic-f132-one-source-schreier-control}
Let
$\mathfrak J=
(\eta_0,\xi_0)$ be
the two-generated
ideal of
Theorem~\ref{thm:dyadic-f132-formal-wild-elimination}.
In the entire completed
genuine $S_3$
residual deformation
ring, all six
Schreier wild exchange
coefficients $b_{ij}$
and all
$M_0\leftrightarrow M_1$
adjoint transport
blocks are divisible
by $\mathfrak J$.
Every two-step
first-diagonal return
defect belongs to
$\mathfrak J^2$,
and every
second-diagonal
return defect
belongs to
$2\mathfrak J^2$.
The three-coset
exchange syzygies
\eqref{eq:dyadic-f132-schreier-zero-sum}
and
\eqref{eq:dyadic-f132-schreier-curvature-zero-sums}
are exact identities
over this \emph{same}
two-generated ideal.

Consequently the
first wild
off-isotypic packet
$b_0$, together
with the complete
marked relation
and the tame
crossed-product
algebra, controls
all higher-order
Schreier exchange
defects at the
coefficient level.
This gives a
strictly smaller
integral source
for constructing
a future
Schreier-to-minimal-Fox
relation-syzygy
comparison.
It is not itself
such a comparison:
no universal
length-two
projective group
resolution or
canonical minimal
rank-five
Demu\v{s}kin basis
is claimed.
\end{corollary}

\begin{proof}
Theorem~\ref{thm:dyadic-f132-formal-wild-elimination}
puts $b_1$ in
$(b_0)$.
The Schreier
packet formula
\eqref{eq:dyadic-f132-schreier-wild-packet}
therefore puts
all six $b_{ij}$
in $\mathfrak J C$.
The full crossed
transport law
\eqref{eq:dyadic-f132-crossed-product-law}
shows that this
property persists
under group-word
multiplication
and inversion,
and the tame
normalizer
preserves the
decomposition.
The closed adjoint
block formulas
\eqref{eq:dyadic-f132-adjoint-four-blocks}
and return
filtration
\eqref{eq:dyadic-f132-return-filtration}
now give the stated
bounds.
The exact
Schreier relations
were proved in
Theorem~\ref{thm:dyadic-f132-schreier-exchange}.
No claim about
identities among
the original
group relators
is made.
\end{proof}

\begin{warning}[Two distinct notions of second-degree syzygy]
\label{warn:dyadic-f132-fox-versus-exchange-syzygy}
The exact identities
$\sum_i b_{ij}=0$
and
$\sum_i\mathscr K_{i\ell}=0$
are \emph{Schreier
coefficient exchange}
relations arising
from the finite
tame orbit and
the quadratic
\'etale crossed
algebra.
They are not
identities among
Roe--Turturean
relators in the
completed group
algebra.
Formal elimination
of the $b_1$
off-isotypic
coordinates
does not remove
the integral
higher-syzygy
problem for the
rank-five maximal
pro-$2$ local
Galois group
$G_{\mathbf Q_2(\sqrt[3]2)}(2)$.
A full finite
Fox chain comparison
must still produce
a verified change
of relation modules
with the correct
degree-two
Fox boundary and
compatibility
with the
index-three
transfer.
The unconditional
all-Sylow perfect
derived atlas
remains valid
independently of
that stronger
construction.
\end{warning}
